% Required Packages: \usepackage{graphicx}, \usepackage{float}

\begin{figure}[htbp]
  \centering
  \begin{subfigure}[b]{0.48\linewidth}
    \centering
    \zimg{rf_fig_20.png}{width=\linewidth,max height=0.4\textheight,keepaspectratio}{fg_0_12qgj}{rf_fig_20.png}
    \caption{(a) shows the behavior of various operational transconductance amplifier (OTA) topologies' DC Gain as the transistors' oxide...}

  \end{subfigure}
  \hfill
  \begin{subfigure}[b]{0.48\linewidth}
    \centering
    \zimg{rf_fig_21.png}{width=\linewidth,max height=0.4\textheight,keepaspectratio}{fg_1_1dbze}{rf_fig_21.png}
    \caption{(a) shows the behavior of various operational transconductance amplifier (OTA) topologies' DC Gain as the transistors' oxide...}

  \end{subfigure}
  \caption{(a) shows the behavior of various operational transconductance amplifier (OTA) topologies' DC Gain as the transistors' oxide thickness (TOX) varies between 2 nm and 5 nm. Gate capacitance and drain current are often significantly impacted by the thickness of the gate oxide. Although it might result in leakage, a thinner oxide often improves gate control. The transconductance (gmn) and output resistance (ROCASN) of the transistors, which determine DC gain (AV=gmn* ROUT), seem to be leveling out or being controlled by other dominating elements in the OTA within this particular range, according to the flat lines \cite{ref7,ref39}.}

  \label{fig:group_y7snq}
\end{figure}
